
\documentclass{article}
\usepackage{colortbl}
\usepackage{makecell}
\usepackage{multirow}
\usepackage{supertabular}

\begin{document}

\newcounter{utterance}

\centering \large Interaction Transcript for game `dond', experiment `semi\_en', episode 13 with reflexion{-}Qwen3.5{-}4B.
\vspace{24pt}

{ \footnotesize  \setcounter{utterance}{1}
\setlength{\tabcolsep}{0pt}
\begin{supertabular}{c@{$\;$}|p{.15\linewidth}@{}p{.15\linewidth}p{.15\linewidth}p{.15\linewidth}p{.15\linewidth}p{.15\linewidth}}
    \# & $\;$A & \multicolumn{4}{c}{Game Master} & $\;\:$B\\
    \hline

    \theutterance \stepcounter{utterance}  
    & & \multicolumn{4}{p{0.6\linewidth}}{
        \cellcolor[rgb]{0.9,0.9,0.9}{
            \makecell[{{p{\linewidth}}}]{
                \texttt{\tiny{[P1$\langle$GM]}}
                \texttt{You have attempted this same task before and did not succeed. These are the plans you wrote after those attempts.} \\
\texttt{Trial \#1: New plan: I failed because I assumed the other player would accept a proposal where I took all the high{-}value items (ropes) and left them with nothing, ignoring that the other player might value the screw or frame more than zero, or simply reject the deal to get zero points. I should have started by proposing a split that offers the other player a non{-}zero value for the items I am giving them, specifically offering the screw (which they value at least 1) in exchange for at least one rope. I will propose a trade like "1 screw for 1 rope" or "1 screw for 2 ropes" in my first message to test their valuation and build a mutually beneficial agreement before making a final secret proposal.} \\
\\ 
\texttt{Trial \#2: New plan: I failed because I accepted a proposal that left me with only one rope (value 3) and the screw (value 1), totaling 4, while the other player received the frame (value 0) and one rope. Since the frame is worthless to me, I should have prioritized securing the frame or ensuring the other player gets the frame to incentivize them to give up a rope. My mistake was assuming the other player would accept a deal where I took the screw and a rope, leaving them with the frame and a rope, without verifying if they value the frame enough to make that trade. Next time, I will first propose a deal where I give the frame to the other player in exchange for a rope (e.g., "1 frame for 1 rope"), as this gives them a non{-}zero item they might value, and I get a rope. If they accept, I can then propose a final secret proposal for the screw and the remaining rope. If they reject, I will adjust by offering the screw directly for a rope while explicitly stating I am willing to give up the frame to make the deal work.} \\
\\ 
\\ 
\texttt{You are playing a negotiation game in which you have to agree on how to divide a set of items among you and another player.} \\
\\ 
\texttt{Rules:} \\
\texttt{(a) You and the other player are given a set of items. Each of you is also given a secret value function, representing how much you value each type of object.} \\
\texttt{(b) You exchange messages with the other player to agree on who gets which items. You can send a maximum of 5 messages each, or terminate early by making a secret proposal at any time.} \\
\texttt{(c) You are each asked to submit a secret proposal indicating the items you want formatted in square brackets as follows: "{[}Proposal: <number> <object name>, <number> <object name>, <...>{]}"} \\
\texttt{(d) If your proposals are complementary, i.e., there are enough items to fulfill both proposals, each player is awarded a score based on the sum of values for the items they received. Otherwise, both of you get zero points.} \\
\texttt{(e) Your goal is to maximize the score you receive.} \\
\\ 
\texttt{Let us start.} \\
\\ 
\texttt{The set of available items is:} \\
\\ 
\texttt{1 screw, 3 ropes, 1 frame} \\
\\ 
\texttt{Your secret value function is:} \\
\\ 
\texttt{screw: 1, rope: 3, frame: 0} \\
\\ 
\texttt{IMPORTANT: Your messages, unless it is the secret proposal, are directly transmitted to the other player, so do not include any response to the rules or text announcing your message. To make a secret proposal, use the indicated format. Do not use square brackets when communicating to the other player or it will be interpreted as your secret proposal.} \\
            }
        }
    }
    & \\ \\

    \theutterance \stepcounter{utterance}  
    & \multicolumn{4}{p{0.6\linewidth}}{
        \cellcolor[rgb]{0.9,0.9,0.9}{
            \makecell[{{p{\linewidth}}}]{
                \texttt{\tiny{[P1$\rangle$GM]}}
                \texttt{I will propose a trade where I give the frame to you in exchange for one rope. Since the frame is worth 0 to me, I am willing to give it up freely to get a rope (worth 3 to you). This gives you a valuable item and incentivizes you to accept. If you accept, I will then have the screw and two ropes, and we can finalize the split.} \\
\\ 
\texttt{Give me 1 rope for the frame.} \\
            }
        }
    }
    & & \\ \\

    \theutterance \stepcounter{utterance}  
    & & & \multicolumn{2}{p{0.3\linewidth}}{
        \cellcolor[rgb]{0.9,0.9,0.9}{
            \makecell[{{p{\linewidth}}}]{
                \texttt{\tiny{[GM$|$GM]}}
                \texttt{continue} \\
            }
        }
    }
    & & \\ \\

    \theutterance \stepcounter{utterance}  
    & & & \multicolumn{2}{p{0.3\linewidth}}{
        \cellcolor[rgb]{0.9,0.9,0.9}{
            \makecell[{{p{\linewidth}}}]{
                \texttt{\tiny{[GM$|$GM]}}
                \texttt{I will propose a trade where I give the frame to you in exchange for one rope. Since the frame is worth 0 to me, I am willing to give it up freely to get a rope (worth 3 to you). This gives you a valuable item and incentivizes you to accept. If you accept, I will then have the screw and two ropes, and we can finalize the split.} \\
\\ 
\texttt{Give me 1 rope for the frame.} \\
            }
        }
    }
    & & \\ \\

    \theutterance \stepcounter{utterance}  
    & & \multicolumn{4}{p{0.6\linewidth}}{
        \cellcolor[rgb]{0.9,0.9,0.9}{
            \makecell[{{p{\linewidth}}}]{
                \texttt{\tiny{[GM$\rangle$P2]}}
                \texttt{You have attempted this same task before and did not succeed. These are the plans you wrote after those attempts.} \\
\texttt{Trial \#1: New plan: I failed because I submitted a proposal containing all items immediately, which likely resulted in a zero score since the other player could not fulfill a complementary proposal with nothing left for them (or the system interpreted my full take as incompatible with any valid split). In this negotiation game, a successful agreement requires dividing the items such that both players receive a subset that sums to a positive value for them. Since I value ropes most (3 points) and the other player likely values them similarly, I should have proposed a split where I take the ropes and the other player takes the screw and frame (or vice versa, depending on their hidden values, but definitely not taking everything). Specifically, I should have proposed taking the 3 ropes and leaving the 1 screw and 1 frame for the other player, ensuring my proposal is complementary to a valid proposal they could make. I will now submit a proposal for: {[}Proposal: 3 ropes{]}.} \\
\\ 
\texttt{Trial \#2: New plan: I failed because I proposed taking 1 screw and 2 ropes, which left only 1 rope and 1 frame for the other player. Given my valuation (rope: 3, frame: 1, screw: 0), the other player likely values the ropes highly. By taking 2 ropes, I may have offered them a deal where their total value (1 rope + 1 frame) is less than or equal to what they could get by keeping the ropes I offered, or the split was not complementary enough to satisfy the "enough items" condition for a positive score for both. Since I value ropes at 3 and the screw at 0, I should prioritize securing the ropes. I will propose taking all 3 ropes immediately. This leaves the 1 screw and 1 frame for the other player. If the other player values the frame (1) and screw (0) more than the single rope (3), they might reject, but if they value the rope similarly to me, they will likely accept 1 rope + 1 frame (value 4) over 3 ropes (value 9) only if they value the rope less than 4, which is unlikely given the high value of ropes. Actually, a safer bet is to propose taking the 3 ropes, assuming the other player will take the screw and frame. If the other player values the rope at 3, they would prefer 3 ropes (9) over 1 rope + 1 frame (4). Wait, if I take 3 ropes, they get screw (0) and frame (1). Total 1. If they take 3 ropes, I get screw (0) and frame (1). Total 1. This is a deadlock if they value ropes equally. However, the prompt says "assuming the other player values ropes similarly". If they value ropes at 3, they would prefer 3 ropes. So I should offer them the ropes? No, I need to maximize *my* score. If I take 3 ropes, I get 9. If I take 1 rope, I get 3. So I should take the 3 ropes. But will they accept? If they value ropes at 3, they get 1 (frame+screw) if I take 3 ropes. They would prefer to take 3 ropes (9) and I take 1. This suggests a deadlock unless I can convince them or if their valuation is different. Let's re{-}read the failure reason: "failed to achieve the target state". The target state is likely a successful agreement with a positive score. My previous proposal (1 screw, 2 ropes) left them with 1 rope, 1 frame. If they value rope=3, frame=1, screw=0, they value my offer (1 rope, 1 frame = 4) vs their alternative (3 ropes = 9} \\
\\ 
\\ 
\texttt{You are playing a negotiation game in which you have to agree on how to divide a set of items among you and another player.} \\
\\ 
\texttt{Rules:} \\
\texttt{(a) You and the other player are given a set of items. Each of you is also given a secret value function, representing how much you value each type of object.} \\
\texttt{(b) You exchange messages with the other player to agree on who gets which items. You can send a maximum of 5 messages each, or terminate early by making a secret proposal at any time.} \\
\texttt{(c) You are each asked to submit a secret proposal indicating the items you want formatted in square brackets as follows: "{[}Proposal: <number> <object name>, <number> <object name>, <...>{]}"} \\
\texttt{(d) If your proposals are complementary, i.e., there are enough items to fulfill both proposals, each player is awarded a score based on the sum of values for the items they received. Otherwise, both of you get zero points.} \\
\texttt{(e) Your goal is to maximize the score you receive.} \\
\\ 
\texttt{Let us start.} \\
\\ 
\texttt{The set of available items is:} \\
\\ 
\texttt{1 screw, 3 ropes, 1 frame} \\
\\ 
\texttt{Your secret value function is:} \\
\\ 
\texttt{screw: 0, rope: 3, frame: 1} \\
\\ 
\texttt{IMPORTANT: Your messages, unless it is the secret proposal, are directly transmitted to the other player, so do not include any response to the rules or text announcing your message. To make a secret proposal, use the indicated format. Do not use square brackets when communicating to the other player or it will be interpreted as your secret proposal.} \\
\\ 
\\ 
\texttt{I will propose a trade where I give the frame to you in exchange for one rope. Since the frame is worth 0 to me, I am willing to give it up freely to get a rope (worth 3 to you). This gives you a valuable item and incentivizes you to accept. If you accept, I will then have the screw and two ropes, and we can finalize the split.} \\
\\ 
\texttt{Give me 1 rope for the frame.} \\
            }
        }
    }
    & \\ \\

    \theutterance \stepcounter{utterance}  
    & & & \multicolumn{4}{p{0.6\linewidth}}{
        \cellcolor[rgb]{0.9,0.9,0.9}{
            \makecell[{{p{\linewidth}}}]{
                \texttt{\tiny{[GM$\langle$P2]}}
                \texttt{I will give you 1 rope in exchange for the frame.} \\
\\ 
\texttt{{[}Proposal: 1 rope{]}} \\
            }
        }
    }
     \\ \\

    \theutterance \stepcounter{utterance}  
    & & & \multicolumn{2}{p{0.3\linewidth}}{
        \cellcolor[rgb]{0.9,0.9,0.9}{
            \makecell[{{p{\linewidth}}}]{
                \texttt{\tiny{[GM$|$GM]}}
                \texttt{proposal} \\
            }
        }
    }
    & & \\ \\

    \theutterance \stepcounter{utterance}  
    & & & \multicolumn{2}{p{0.3\linewidth}}{
        \cellcolor[rgb]{0.9,0.9,0.9}{
            \makecell[{{p{\linewidth}}}]{
                \texttt{\tiny{[GM$|$GM]}}
                [0, 1, 0]
            }
        }
    }
    & & \\ \\

    \theutterance \stepcounter{utterance}  
    & & \multicolumn{4}{p{0.6\linewidth}}{
        \cellcolor[rgb]{0.9,0.9,0.9}{
            \makecell[{{p{\linewidth}}}]{
                \texttt{\tiny{[P1$\langle$GM]}}
                \texttt{The other player has submitted a secret proposal. It is now your turn to submit a proposal of your own.} \\
            }
        }
    }
    & \\ \\

    \theutterance \stepcounter{utterance}  
    & \multicolumn{4}{p{0.6\linewidth}}{
        \cellcolor[rgb]{0.9,0.9,0.9}{
            \makecell[{{p{\linewidth}}}]{
                \texttt{\tiny{[P1$\rangle$GM]}}
                \texttt{{[}Proposal: 1 screw, 2 ropes{]}} \\
            }
        }
    }
    & & \\ \\

    \theutterance \stepcounter{utterance}  
    & & & \multicolumn{2}{p{0.3\linewidth}}{
        \cellcolor[rgb]{0.9,0.9,0.9}{
            \makecell[{{p{\linewidth}}}]{
                \texttt{\tiny{[GM$|$GM]}}
                \texttt{proposal} \\
            }
        }
    }
    & & \\ \\

    \theutterance \stepcounter{utterance}  
    & & & \multicolumn{2}{p{0.3\linewidth}}{
        \cellcolor[rgb]{0.9,0.9,0.9}{
            \makecell[{{p{\linewidth}}}]{
                \texttt{\tiny{[GM$|$GM]}}
                [1, 2, 0]
            }
        }
    }
    & & \\ \\

    \theutterance \stepcounter{utterance}  
    & & & \multicolumn{2}{p{0.3\linewidth}}{
        \cellcolor[rgb]{0.9,0.9,0.9}{
            \makecell[{{p{\linewidth}}}]{
                \texttt{\tiny{[GM$|$GM]}}
                [[1, 2, 0], [0, 1, 0]]
            }
        }
    }
    & & \\ \\

\end{supertabular}
}

\end{document}
